\section{Limitations and Future Directions}
\label{sec:limitations-fur-dir}

\paragraph{Cost and Runtime.}
Although \ourmethod improves success rates on long-horizon shared-artifact tasks, it introduces non-trivial coordination overhead. In our experiments, multi-agent execution consistently incurs higher API cost than single-agent baselines, and wall-clock runtime is not substantially reduced despite parallel execution. This reflects a fundamental trade-off: structured isolation, integration, and verification improve stability, but require additional communication rounds, merge operations, and test executions. In particular, while engineers operate concurrently, integration remains sequential and test-gated, limiting end-to-end acceleration. Prior analyses of multi-agent systems have similarly noted that coordination complexity can offset gains from specialization and parallelism when not carefully optimized~\citep{oreilly2024multiagent}. For the long-horizon shared-artifact tasks we study, however, such coordination may still be necessary, since simply extending single-agent execution does not reliably achieve comparable gains. Therefore, promising next steps include improving scheduling efficiency, reducing redundant verification cycles, and learning when to merge or prune intermediate states. Optimizing the cost–performance frontier of structured multi-agent execution remains an important area for future work. 

% In particular, while engineers operate concurrently, integration remains sequential and test-gated, limiting end-to-end acceleration. Prior analyses of multi-agent systems have similarly noted that coordination complexity can offset gains from specialization and parallelism when not carefully optimized~\citep{oreilly2024multiagent}. 

% Improving scheduling efficiency, reducing redundant verification cycles, and learning when to merge or prune intermediate states are promising directions for reducing this overhead. Optimizing the cost–performance frontier of structured multi-agent execution remains an important area for future work. 


\paragraph{Isolated Task Delegation Capabilities of Agents.}
A second limitation lies in the central manager’s task decomposition and delegation capability. In the current implementation, task assignment relies primarily on prompt engineering heuristics rather than learned delegation policies. While our results indicate that architectural isolation and integration are the dominant factors for stability, weak or suboptimal task decomposition can still reduce overall effectiveness. Existing analyses of multi-agent systems identify imprecise task handoffs and underspecified subgoals as major sources of coordination failure~\citep{galileo2026failures}. 
% \gncomment{It shouldn't be just ``Labs'' in the reference, use the full name.}. 
Our findings align with this observation: when delegation is coarse-grained or misaligned with dependency structure, engineers may produce locally correct outputs that are globally inefficient to integrate. Future work may explore reinforcement learning–based delegation policies, dependency-aware planning modules, or adaptive subtask refinement strategies that improve alignment between global objectives and isolated execution. Strengthening delegation capability would allow the architectural benefits of isolation and structured integration to scale more reliably.

\paragraph{Generalization Beyond Software Engineering Tasks.}
Finally, our evaluation focuses on software engineering benchmarks, which provide a natural testbed for structured multi-agent execution due to explicit workspace boundaries, version control infrastructure, and executable test suites. These properties make software development uniquely suitable for studying isolation, integration, and dependency-aware coordination. However, not all long-horizon shared-artifact tasks possess such clearly defined boundaries or objective verification mechanisms. Extending \ourmethod to non-coding domains—such as document synthesis, research planning, or multimodal artifact construction—will require adapting isolation mechanisms and designing alternative forms of integration and validation. Evaluating the framework in such settings is necessary to determine whether the architectural principles demonstrated here generalize beyond SWE-specific workflows.